\documentclass[12pt,a4paper]{report}

% --- Pakiety podstawowe ---
\usepackage[utf8]{inputenc}
\usepackage[T1]{fontenc}
\usepackage[polish]{babel}
\usepackage{amsmath}
\usepackage{amsfonts}
\usepackage{amssymb}
\usepackage{graphicx}
\usepackage{booktabs}
\usepackage{hyperref}
\usepackage{listings}
\usepackage{geometry}

\geometry{margin=2.5cm}

% --- Metadane ---
\title{Uczenie silnika szachowego na partiach o rosnącym Elo}
\author{Adam Zieliński}
\date{\today}

\begin{document}

% --- Strona tytułowa (uproszczona) ---
\begin{titlepage}
    \centering
    \vspace*{1cm}
    {\Large Uniwersytet Mikołaja Kopernika w Toruniu\\ Wydział Matematyki i Informatyki \par}
    \vspace{3cm}
    {\Huge \textbf{Uczenie silnika szachowego na partiach o rosnącym Elo} \par}
    \vspace{2cm}
    {\Large Praca magisterska \par}
    \vspace{2cm}
    {\Large Autor: Adam Zieliński \par}
    {\Large Promotor: dr hab. Łukasz Mikulski, prof. UMK \par}
    \vfill
    {\Large Toruń, 2026 \par}
\end{titlepage}

\tableofcontents

\chapter*{Streszczenie}

Niniejsza praca bada wpływ uczenia progresywnego (\textit{curriculum learning}) na proces trenowania głębokich funkcji ewaluacji w silnikach szachowych opartych na architekturze NNUE (\textit{Efficiently Updatable Neural Network}). Głównym celem badawczym było zweryfikowanie hipotezy, czy sekwencyjne douczanie modelu na partiach pochodzących z coraz wyższych przedziałów rankingowych (Elo) pozytywnie wpływa na zbieżność procesu optymalizacji, dynamikę funkcji straty oraz końcową siłę gry w porównaniu do podejścia klasycznego (trening na zbiorze o wymieszanym rozkładzie rankingowym).

Projekt opiera się na koncepcji rozwijanej od etapu pracy inżynierskiej, w ramach której dążono do stworzenia adaptacyjnego asystenta treningowego, dostosowującego poziom trudności do umiejętności gracza w celu wywołania tzw. stanu optymalnego zaangażowania (\textit{flow}). W toku wstępnych eksperymentów odrzucono podejście oparte na architekturze Leela Chess Zero ze względu na ograniczenia infrastrukturalno-dokumentacyjne, decydując się na środowisko \texttt{nnue-pytorch} oraz zmodyfikowany silnik Stockfish (z funkcją cech \textit{HalfKAv2\_hm\string^}).

Zbiorami danych wejściowych była oficjalna baza danych serwisu Lichess z marca 2026 roku (ponad 90 milionów pojedynków). Przy użyciu autorskiego potoku przetwarzania (skrypty w języku Rust oraz narzędzie \texttt{pgn-extract}) surowe partie PGN przefiltrowano pod kątem minimalnej długości wariantów (\texttt{--minply 40}), poprawności zakończeń (\texttt{--nobadresults}) oraz wyeliminowania asymetrii rankingowej. Przefiltrowane dane podzielono na rozłączne koszyki Elo (po 2 miliony gier w każdym), a następnie przetransformowano je do binarnego formatu \texttt{.binpack}.

W celach weryfikacyjnych przeprowadzono eksperymenty porównawcze obejmujące trzy próby:
\begin{enumerate}
    \item \textbf{Badawczą (progresywną):} uczenie realizowane kolejno od najniższych do najwyższych koszyków Elo.
    \item \textbf{Kontrolną:} uczenie na losowo wymieszanym zbiorze danych z pełnego zakresu Elo.
    \item \textbf{Odwróconą (\textit{reverse curriculum}):} uczenie rozpoczynane od partii o wysokim Elo, a kończone na partiach graczy początkujących.
\end{enumerate}

Analiza statystyczna krzywych funkcji straty (\textit{loss}) wykazała istotne różnice w dynamice uczenia pomiędzy modelem kontrolnym a próbami sekwencyjnymi. W próbie badawczej odnotowano aż 2,5-krotnie wyższe całkowite odchylenie standardowe błędu ($\text{Total Loss Std} = 0{,}006142$ vs $0{,}002484$) oraz o 81\% wyższą zmienność szumu względem trendu ($\text{Noise Volatility} = 0{,}002060$ vs $0{,}001137$). Sytuację tę wywołały gwałtowne, punktowe skoki funkcji straty w momentach przełączania zasobników PGN.

Najważniejszym odkryciem pracy jest symetria zachowania sieci w próbie progresywnej i odwróconej. Wskaźniki całkowitej zmienności błędu oraz liczby nagłych skoków ($>3\sigma$) okazały się niemal identyczne w obu tych wariantach (różnice poniżej 4\%). Wynik ten jednoznacznie obala hipotezę, jakoby partie graczy o wyższym rankingu Elo stanowiły dla sieci problem o „wyższym stopniu trudności optymalizacyjnej”. Obserwowane destabilizacje są wyłącznie konsekwencją dyskontynuacji rozkładu danych uczących (\textit{Co-Domain Shift} / \textit{Distribution Shift}). 

Praca formułuje powiązane wnioski praktyczne dla projektowania systemów uczących się na bazach gier ludzkich: sekwencyjna podmiana zbiorów danych bez zastosowania mechanizmów ciągłego powtarzania danych (\textit{Data Replay}) prowadzi do zaburzeń procesu optymalizacji wag, niezależnie od kierunku zmian jakościowych w danych.

\chapter{Wstęp}

\section{Motywacja}
% Na podstawie README_AI.md: Rozwinięcie pracy inżynierskiej, asystent uczący się z graczem.

Pierwszym celem jaki chciałem osiągnąć było stworzenie partnera do trenowania gry w szachy, w taki sposób aby zawsze był na poziomie gracza, potencjalnie lekko silniejszym, aby gracz podczas każdej rozgrywki mógł doświadczać tzw. stanu flow, gdzie nie czuje nudy przez zbyt słabego przeciwnika ani frustracji gdy jest zbyt silny. Aby dopasować poziom trudności silnika, uznałem że dobrym rozwiązaniem byłoby douczanie sieci za każdym razem gdy procent wygranych gracza przekroczy pewną wartość. Wtedy danymi uczącymi byłyby rozgrywki lekko wyprzedzające aktualne ELO gracza. Aby zbadać czy taka implementacja spełniałaby moje oczekiwania, uznałem że wpierw przeprowadzę eksperyment aby zbadać jak będzie zachowywać się silnik uczony grami o coraz to lepszej jakości. 

\section{Cel pracy}
% Badanie zachowania silnika NN podczas uczenia na danych o rosnącym poziomie umiejętności.

\section{Hipotezy badawcze}
\begin{itemize}
    \item \textbf{H1: Szybsza nauka cech.} Model szybciej odnajdzie wzorce na prostszych danych.
    \item \textbf{H2: Ludzki styl gry.} Model będzie unikał losowości typowej dla silników uczonych na "czystej taktyce".
    \item \textbf{H3: Korelacja Elo.} Zależność między rankingiem zbioru treningowego a siłą gry modelu.
\end{itemize}

\chapter{Przegląd technologii}

\section{Silnik Stockfish i architektura NNUE}
% Opis Stockfish NNUE, nnue-pytorch.
Stockfish to jeden z najsilniejszych na świecie, otwartoźródłowych silników szachowych, rozwijany społecznościowo od 2008 roku. Tradycyjnie opierał się na algorytmie przeszukiwania drzewa wariantów Alpha-Beta uzupełnionym o rozbudowane, ręcznie projektowane funkcje ewaluacji (Hand-Crafted Evaluation - HCE), które oceniały pozycję na podstawie tysięcy reguł czysto taktycznych i pozycyjnych.
Przełomem w rozwoju silnika było wprowadzenie w 2020 roku architektury NNUE (Efficiently Updatable Neural Network), pierwotnie zaprojektowanej przez Yu Nasu dla shōgi. NNUE to płytka, wolna od zwijania (fully connected) sieć neuronowa zintegrowana bezpośrednio z silnikiem przeszukującym, która zastąpiła tradycyjną funkcję HCE. Główną zaletą tej architektury jest jej dopasowanie do specyfiki gry w szachy - przy wykonywaniu pojedynczego ruchu na planszy większość cech pozycji nie ulega zmianie, co pozwala na przyrostowe (inkrementalne) aktualizowanie wartości aktywacji pierwszej warstwy sieci zamiast przeliczania całej pozycji od zera. W praktyce sieć neuronowa NNUE działa bezpośrednio na procesorze (CPU) z wykorzystaniem rozkazów SIMD (np. AVX2, AVX-512), osiągając przepustowość milionów ewaluacji na sekundę. W wersjach takich jak \textit{HalfKAv2\_hm\string^}) stan planszy reprezentowany jest jako rzadki wektor cech opisujący relacje pomiędzy pozycją króla a położeniem pozostałych bierek, co pozwala modelowi na precyzyjną ocenę dynamiki pozycji przy zachowaniu wyjątkowo wysokiej wydajności obliczeniowej.

\section{Baza danych Lichess}
Lichess.org jest jedną z największych na świecie otwartych platform szachowych, regularnie udostępniającą publicznie kompletne, zaanonimowane bazy danych partii rozegranych przez użytkowników serwisu. Zbiór ten stanowi idealne źródło danych treningowych na potrzeby eksperymentów z uczeniem maszynowym ze względu na swoją reprezentatywność, wolumen oraz szczegółowość metadanych.
Każda partia w bazie zapisana jest w formacie PGN (Portable Game Notation) i zawiera m.in. dokładne rankingi Elo obu graczy, kontrolę czasu (od partii błyskawicznych po klasyczne) oraz uzgodniony wynik. Różnorodność poziomów umiejętności graczy - od początkujących (Elo poniżej 1000) po arcymistrzów (Elo powyżej 2700) - jest kluczowym elementem z punktu widzenia koncepcji uczenia progresywnego (curriculum learning). Umożliwia to precyzyjne selekcjonowanie partii i podział danych na spójne podbiory reprezentujące kolejne etapy rozwoju szachowego, co stanowi fundament metodologiczny niniejszej pracy.

\section{Narzędzia pomocnicze}
% pgn-extract, cutechess, Ordo.
Proces przetwarzania danych, przeprowadzenia treningu oraz automatyzacji turniejów ewaluacyjnych wymagał wykorzystania wyspecjalizowanego zestawu narzędzi open-source:
\begin{itemize}
    \item \textbf{pgn-extract} zaawansowane narzędzie wiersza poleceń przeznaczone do przeszukiwania, filtrowania i parsowania plików PGN. W ramach przeprowadzonego eksperymentu posłużyło do selekcji partii według kryteriów takich jak minimalna liczba posunięć (filtracja zagrań zbyt krótkich lub porzuconych), poziom Elo graczy oraz odrzucenie partii zakończonych nietypowo (np. z powodu problemów technicznych z połączeniem).
    \item \textbf{Cutechess} elastyczny interfejs i narzędzie wiersza poleceń służące do automatycznego przeprowadzania turniejów pomiędzy silnikami szachowymi operującymi na protokole UCI (Universal Chess Interface). Pozwoliło na obiektywną weryfikację siły gry poszczególnych iteracji modelu w kontrolowanych warunkach (z jednakową kontrolą czasu i zdefiniowanymi pozycjami otwarć).
    \item \textbf{Ordo} narzędzie statystyczne dedykowane do wyznaczania rankingów Elo na podstawie wyników bezpośrednich pojedynków pomiędzy silnikami. Pozwala na precyzyjne obliczenie siły gry modelu wraz z przedziałami ufności na podstawie rozegranych turniejów w Cutechess.
\end{itemize}

\chapter{Przygotowanie danych}

\section{Pozyskiwanie i filtrowanie danych (PGN)}
% Opis użycia pgn-extract, parametry minply, nobadresults.
Proces przygotowania zbioru treningowego rozpoczęto od pobrania surowych archiwów PGN z otwartego repozytorium bazy danych Lichess. Ze względu na fakt, że surowa baza zawiera miliony partii o bardzo zróżnicowanej jakości, konieczne było przeprowadzenie rygorystycznej filtracji w celu odrzucenia rekordów niepełnych, błędnych lub zaburzających proces uczenia.
Do wstępnej obróbki plików PGN wykorzystano program pgn-extract. Zastosowano następujące reguły i parametry filtrowania:
\begin{itemize}
    \item \textbf{Minimalna długość partii (--minply / -s)} Odrzucono partie krótsze niż określony próg posunięć (np. poniżej 10 półruchów). Pozwoliło to wyeliminować szybkie remisy polubowne, natychmiastowe poddania w debiucie oraz partie zerwane na bardzo wczesnym etapie.
    \item \textbf{textbfOdrzucenie nietypowych zakończeń (--nobadresults)} Przefiltrowano partie zakończone z przyczyn technicznych (porzucenie połączenia, awaria serwera, błędy protokołu), pozostawiając jedynie pojedynki rozstrzygnięte w sposób naturalny (mat, poddanie, pat, przekroczenie czasu lub wyczerpanie materiału).
    \item \textbf{Jednolity poziom graczy} Odrzucono partie z dużą dysproporcją rankingową pomiędzy zawodnikami, aby uniknąć sytuacji, w której wysoce asymetryczna jakość zagrań zaburzałaby ocenę cech pozycji.
\end{itemize}

\section{Podział na koszyki Elo}
% Algorytm podziału gier na przedziały rankingowe.


\section{Transformacja do formatu binpack}
% pgn -> plain -> binpack.
Prawidłowy i wydajny trening architektury NNUE za pomocą biblioteki nnue-pytorch wymaga przekształcenia tekstu PGN do dedykowanego, skompresowanego formatu binarnego .binpack. Proces ten realizowany jest w dwuetapowym potoku (pipeline):

Konwersja PGN do formatu pośredniego (plain text): Zapis PGN parsowany jest do uproszczonej postaci tekstowej, w której każda linia reprezentuje stan planszy (w notacji FEN) połączony z wykonanym ruchem, wynikiem partii oraz ewentualną oceną ewaluacyjną silnika.

Kompresja do formatu .binpack: Za pomocą wbudowanych narzędzi konwertujących nnue-pytorch, plik pośredni ulega binarnej kompresji. Format .binpack rezerwuje niewielką liczbę bitów na reprezentację roszad, praw do bicia w przelocie oraz rzadkich wektorów pozycji bierek.
Transformacja ta redukuje rozmiar danych na dysku o ponad 80\% w porównaniu do surowego PGN oraz drastycznie zwiększa przepustowość odczytu (I/O) podczas treningu na GPU, eliminując wąskie gardło CPU podczas rozpakowywania rekordu partii w czasie rzeczywistym.

\chapter{Proces uczenia}

\section{Architektura eksperymentu}

Eksperyment badawczy zaprojektowano w celu dokładnego porównania dynamiki uczenia oraz końcowych właściwości modelu poddawanego uczeniu progresywnemu (\textit{curriculum learning}) względem podejścia klasycznego (kontrolnego) oraz wariantu odwróconego (\textit{reverse curriculum}). Przepływ danych oraz poszczególne etapy przetwarzania i ewaluacji przedstawiono na schemacie (Rysunek~\ref{fig:workflow}).

\begin{figure}[htbp]
    \centering
    % Miejsce na wygenerowany wykres z diagramu Draw.io
    \caption{Schemat przepływu danych oraz architektury eksperymentu badawczego.}
    \label{fig:workflow}
\end{figure}

Cały proces eksperymentalny składa się z następujących kroków:
\begin{enumerate}
    \item \textbf{Przygotowanie i przefiltrowanie bazy danych PGN:} Surowe partie pobrane z serwisu Lichess poddawane są czyszczeniu przy użyciu programu \texttt{pgn-extract} w celu odrzucenia zagrań zbyt krótkich oraz błędnych wyników.
    \item \textbf{Podział na koszyki Elo oraz transformacja danych:} W wariancie badawczym przefiltrowane partie rozdzielane są na osobne zbiorniki reprezentujące rosnące przedziały rankingowe (po 2~miliony partii na koszyk). W przypadku próby kontrolnej partie z różnych przedziałów są losowo wymieszane w ramach jednego zbioru. Następnie dane są ewaluowane i konwertowane do binarnego formatu \texttt{.binpack}.
    \item \textbf{Iteracyjny trening sieci neuronowej:} Uczenie odbywa się przy użyciu biblioteki \texttt{nnue-pytorch}. W wariancie badawczym model po osiągnięciu założonej liczby epok (100 epok na koszyk) przechodzi do kolejnego zbiornika PGN (pętla \textit{Zmiana koszyka}), kontynuując naukę z punktu kontrolnego (\textit{checkpoint}) poprzedniego etapu.
    \item \textbf{Obliczenie metryk pośrednich i ewaluacja:} Po zakończeniu określonych epok następuje serializacja wag sieci do plików \texttt{.nnue}, budowa dedykowanych wersji silnika Stockfish oraz automatyczne przeprowadzanie turniejów testowych (za pomocą narzędzi \texttt{cutechess} i \texttt{Ordo}).
    \item \textbf{Eksploracja danych i wnioskowanie:} Zgromadzone krzywe funkcji straty (\textit{loss curves}) oraz metryki rankingowe poddawane są analizie statystycznej w celu weryfikacji postawionych hipotez.
\end{enumerate}

\section{Parametry treningu}

Trening sieci realizowano z wykorzystaniem frameworka \texttt{nnue-pytorch} (opartego na \texttt{PyTorch Lightning}). Zastosowano spójny zestaw hiperparametrów dla wszystkich analizowanych wariantów eksperymentalnych:

\begin{itemize}
    \item \textbf{Reprezentacja cech (\textit{Feature Set}):} Zastosowano układ \textit{HalfKAv2\_hm\string^}, będący standardem w nowoczesnych wersjach silnika Stockfish. Koduje on względne pozycje figur na planszy w relacji do położenia obu królów z uwzględnieniem symetrii planszy.
    \item \textbf{Wielkość wsadu (\textit{Batch Size}):} Trening prowadzono z rozmiarem wsadu równym $16\,384$ pozycji. Wybór tak dużego wsadu zapewnia wysoką stabilność estymacji gradientu na kartach graficznych o dużej pamięci VRAM oraz optymalne wykorzystanie mocy obliczeniowej GPU.
    \item \textbf{Współczynnik straty (\textit{Lambda}):} Parametr $\lambda = 0{,}3$ reguluje wagę pomiędzy oceną wartości pozycji (\textit{eval}) wynikającej ze statycznej oceny silnika a rzeczywistym wynikiem partii (\textit{game outcome}).
    \item \textbf{Liczba epok i harmonogram:} Dla każdego koszyka rankingowego zaplanowano po 100 epok uczenia (co sumarycznie dla 6 koszyków przekłada się na 600 epok i około 50--80 godzin ciągłych obliczeń). Zapis stanu sieci (\textit{checkpoint}) wykonywany był co 10 epok (\texttt{--network-save-period 10}).
\end{itemize}

Przykładowe uruchomienie procesu uczenia wewnątrz środowiska przedstawia poniższy listing:

\begin{lstlisting}[language=bash, caption=Polecenie uruchamiające proces uczenia na pierwszym koszyku Elo]
python train.py \
    /data/input/binpack/clean_1200_under_2mil.binpack \
    --validation-datasets /data/input/binpack/clean_1200_under_200k.binpack \
    --default-root-dir /data/trained_nets/ \
    --gpus 0 \
    --threads 4 \
    --batch-size 16384 \
    --lambda 0.3 \
    --max-epochs 100 \
    --features "HalfKAv2_hm^" \
    --network-save-period 10
\end{lstlisting}

\section{Infrastruktura obliczeniowa}

Ze względu na skomplikowane zależności bibliotek (\texttt{PyTorch}, \texttt{CUDA}/\texttt{ROCm}, \texttt{PyTorch Lightning}) oraz wymóg odizolowania środowiska uruchomieniowego, cały proces uczenia oraz serializacji modeli przeprowadzono wewnątrz kontenera \textbf{Docker}.

Środowisko sprzętowo-programowe wykorzystane w eksperymencie obejmowało:
\begin{itemize}
    \item \textbf{Jednostkę obliczeniową (GPU):} Karta graficzna AMD Radeon RX 6800 XT z 16 GB pamięci VRAM, wykorzystująca przyspieszenie sprzętowe w środowisku sterowników ROCm (przekazywanie urządzeń \texttt{/dev/kfd} oraz \texttt{/dev/dri} do kontenera).
    \item \textbf{Konfigurację kontenera:} Kontener uruchamiano z prawami dostępu do grup systemowych \texttt{render} oraz \texttt{video}, z nieskończonym limitem blokowania pamięci (\texttt{--ulimit memlock=-1}) oraz zwiększonym rozmiarem stosu, co wyeliminowało wąskie gardła podczas przesyłania dużych zbiorów \texttt{.binpack} z dysku do pamięci GPU.
    \item \textbf{Środowisko wykonawcze (Python):} Środowisko wewnątrz kontenera zostało dostosowane pod kątem spójności wersji pakietów (m.in. wymuszenie wersji \texttt{numpy < 2.5} oraz odpowiednich wersji \texttt{torch} kompatybilnych z architekturą sterowników).
\end{itemize}

\chapter{Ewaluacja i wyniki}
\section{Metryki oceny}
\begin{itemize}
    \item Ranking Elo (turnieje Cutechess).
    \item Zgodność z ruchami ludzkimi (Human-like Index).
    \item Krzywe funkcji straty (Loss Curve).
\end{itemize}
\section{Wyniki eksperymentu}
% Analiza porównawcza epok.
\section{Analiza skoków rozwojowych}
% Czy zmiana koszyka Elo powoduje gwałtowny wzrost siły gry?

\chapter{Wnioski}
\section{Weryfikacja hipotez}
\section{Zastosowania praktyczne}
% Interaktywny asystent szachowy.
\section{Kierunki dalszych badań}

\begin{thebibliography}{99}
    \bibitem{stockfish} Stockfish Team, \textit{Stockfish: A powerful open-source chess engine}, \url{https://github.com/official-stockfish/Stockfish}.
    \bibitem{nnuepytorch} Official Stockfish, \textit{nnue-pytorch}, \url{https://github.com/official-stockfish/nnue-pytorch}.
    \bibitem{lichess} Lichess Database, \url{https://database.lichess.org}.
\end{thebibliography}

\listoffigures
\listoftables

\chapter*{Dodatek A: Skrypty automatyzujące}
\begin{lstlisting}[language=bash, caption=Przykładowe uruchomienie treningu]
python train.py \
    /data/input/binpack/clean_1200_under_2mil.binpack \
    --validation-datasets /data/input/binpack/clean_1200_under_200k_test.binpack \
    --gpus 0 \
    --lambda 0.3 \
    --features "HalfKAv2_hm^"
\end{lstlisting}

\end{document}